\section{Problemas}

\question{
Dadas as matrizes
\[
A =
\begin{bmatrix}
1 & 2 \\
1 & 0
\end{bmatrix}
\quad \text{e} \quad
B =
\begin{bmatrix}
3 & -1 \\
0 & 1
\end{bmatrix},
\]
calcule:
\begin{enumerate}
    \item[(a)] $\det A + \det B$
    \item[(b)] $\det (A+B)$
\end{enumerate}
}
\answer{
  \item[(a)] $\det A + \det B$
  \begin{equation}
  \begin{aligned}
    \det A &= 
    \begin{vmatrix}
    1 & 2 \\
    1 & 0
    \end{vmatrix} \\
    &= (1)(0) - (2)(1) \\
    &= -2
  \end{aligned}
  \end{equation}

  \begin{equation}
  \begin{aligned}
    \det B &= 
    \begin{vmatrix}
    3 & -1 \\
    0 & 1
    \end{vmatrix} \\
    &= (3)(1) - (-1)(0) \\
    &= 3
  \end{aligned}
  \end{equation}

  \begin{equation}
  \begin{aligned}
    \det A + \det B &= -2 + 3 \\
                    &= 1
  \end{aligned}
  \end{equation}

  \item[(b)] $\det (A+B)$
  \begin{equation}
  \begin{aligned}
    A+B &=
    \begin{bmatrix}
    1+3 & 2+(-1) \\
    1+0 & 0+1
    \end{bmatrix} \\
    &=
    \begin{bmatrix}
    4 & 1 \\
    1 & 1
    \end{bmatrix}
  \end{aligned}
  \end{equation}

  \begin{equation}
  \begin{aligned}
    \det(A+B) &=
    \begin{vmatrix}
    4 & 1 \\
    1 & 1
    \end{vmatrix} \\
    &= (4)(1) - (1)(1) \\
    &= 3
  \end{aligned}
  \end{equation}
}

\newpage
\question{
Calcule:
\begin{enumerate}
    \item[(a)]
    \[
    A = \det
    \begin{bmatrix}
    3 & -1 & 5 \\
    0 & 2 & 0 \\
    0 & 10 & 0
    \end{bmatrix}
    \]

    \item[(b)]
    \[
    A = \det
    \begin{bmatrix}
    3 & 0 & 0 \\
    19 & 18 & 0 \\
    -6 & \pi & -5
    \end{bmatrix}
    \]

    \item[(c)]
    \[
    A = \det
    \begin{bmatrix}
    -1 & 3 & 2 \\
    3 & -1 & 1 \\
    2 & 1 & -1
    \end{bmatrix}
    \]

    \item[(d)]
    \[
    A = \det
    \begin{bmatrix}
    2 & 3 & 1 \\
    5 & 3 & 1 \\
    0 & 1 & 2
    \end{bmatrix}
    \]
\end{enumerate}
}
\answer{
  \item[(a)]
  \begin{equation}
  \begin{aligned}
    A &= \det
    \begin{bmatrix}
    3 & -1 & 5 \\
    0 & 2 & 0 \\
    0 & 10 & 0
    \end{bmatrix} \\[1em]
    &= 3
    \begin{vmatrix}
    2 & 0 \\
    10 & 0
    \end{vmatrix}
    - (-1)
    \begin{vmatrix}
    0 & 0 \\
    0 & 0
    \end{vmatrix}
    + 5
    \begin{vmatrix}
    0 & 2 \\
    0 & 10
    \end{vmatrix} \\[1em]
    &= 3(0) + 1(0) + 5(0) \\[1em]
    &= 0
  \end{aligned}
  \end{equation}

  \item[(b)]
  \begin{equation}
  \begin{aligned}
    A &= \det
    \begin{bmatrix}
    3 & 0 & 0 \\
    19 & 18 & 0 \\
    -6 & \pi & -5
    \end{bmatrix}
  \end{aligned}
  \end{equation}

  Como a matriz é triangular inferior:
  \begin{equation}
  \begin{aligned}
    A &= (3)(18)(-5) \\
      &= -270
  \end{aligned}
  \end{equation}

  \item[(c)]
  \begin{equation}
  \begin{aligned}
    A &= \det
    \begin{bmatrix}
    -1 & 3 & 2 \\
    3 & -1 & 1 \\
    2 & 1 & -1
    \end{bmatrix} \\[1em]
    &= (-1)
    \begin{vmatrix}
    -1 & 1 \\
    1 & -1
    \end{vmatrix}
    - 3
    \begin{vmatrix}
    3 & 1 \\
    2 & -1
    \end{vmatrix}
    + 2
    \begin{vmatrix}
    3 & -1 \\
    2 & 1
    \end{vmatrix} \\[1em]
    &= (-1)[(-1)(-1)-(1)(1)]
    - 3[(3)(-1)-(1)(2)]
    + 2[(3)(1)-(-1)(2)] \\[1em]
    &= (-1)(1-1)
    - 3(-3-2)
    + 2(3+2) \\[1em]
    &= 0 + 15 + 10 \\[1em]
    &= 25
  \end{aligned}
  \end{equation}

  \item[(d)]
  \begin{equation}
  \begin{aligned}
    A &= \det
    \begin{bmatrix}
    2 & 3 & 1 \\
    5 & 3 & 1 \\
    0 & 1 & 2
    \end{bmatrix} \\[1em]
    &= 2
    \begin{vmatrix}
    3 & 1 \\
    1 & 2
    \end{vmatrix}
    - 3
    \begin{vmatrix}
    5 & 1 \\
    0 & 2
    \end{vmatrix}
    + 1
    \begin{vmatrix}
    5 & 3 \\
    0 & 1
    \end{vmatrix} \\[1em]
    &= 2[(3)(2)-(1)(1)]
    - 3[(5)(2)-(1)(0)]
    + 1[(5)(1)-(3)(0)] \\[1em]
    &= 2(5)-3(10)+5 \\[1em]
    &= 10-30+5 \\[1em]
    &= -15
  \end{aligned}
  \end{equation}
}

\question{
Dadas as matrizes
\[
A =
\begin{bmatrix}
a & c+2a \\
b & d+2b
\end{bmatrix}
\quad \text{e} \quad
B =
\begin{bmatrix}
a & c \\
b & d
\end{bmatrix},
\]
mostre que $\det A = \det B$.
}
\answer{
  \begin{equation}
  \begin{aligned}
    \det A &=
    \begin{vmatrix}
    a & c+2a \\
    b & d+2b
    \end{vmatrix} \\[1em]
    &= a(d+2b) - b(c+2a) \\[1em]
    &= ad + 2ab - bc - 2ab \\[1em]
    &= ad - bc
  \end{aligned}
  \end{equation}

  \begin{equation}
  \begin{aligned}
    \det B &=
    \begin{vmatrix}
    a & c \\
    b & d
    \end{vmatrix} \\[1em]
    &= ad - bc
  \end{aligned}
  \end{equation}

  Portanto:
  \begin{equation}
  \begin{aligned}
    \det A &= \det B
  \end{aligned}
  \end{equation}
}

\newpage
\question{
Mostre que
\[
\begin{vmatrix}
a & b & c \\
d & e & f \\
d & b+e & c+f
\end{vmatrix}
=
-a
\begin{vmatrix}
b & c \\
e & f
\end{vmatrix}.
\]
}
\answer{
  \begin{equation}
  \begin{aligned}
  &
  \begin{vmatrix}
  a & b & c \\
  d & e & f \\
  d & b+e & c+f
  \end{vmatrix}
  \\[1em]
  &= 
  a
  \begin{vmatrix}
  e & f \\
  b+e & c+f
  \end{vmatrix}
  -
  b
  \begin{vmatrix}
  d & f \\
  d & c+f
  \end{vmatrix}
  +
  c
  \begin{vmatrix}
  d & e \\
  d & b+e
  \end{vmatrix}
  \\[1em]
  &=
  a\,[e(c+f)-f(b+e)]
  -
  b\,[d(c+f)-fd]
  +
  c\,[d(b+e)-ed]
  \\[1em]
  &=
  a(ec+ef-fb-fe)
  -
  b(dc)
  +
  c(db)
  \\[1em]
  &=
  a(ec-fb) - bdc + cdb
  \\[1em]
  &=
  a(ec-fb)
  \\[1em]
  &=
  -a(bf-ce)
  \\[1em]
  &=
  -a
  \begin{vmatrix}
  b & c \\
  e & f
  \end{vmatrix}
  \end{aligned}
  \end{equation}
}

\newpage
\question{
O determinante da matriz
\[
A =
\begin{bmatrix}
1 & 0 & -1 \\
0 & x & 0 \\
x & 0 & -1
\end{bmatrix}
\]
é positivo em qual intervalo de $x \in \mathbb{R}$?
}
\answer{
  \begin{equation}
  \begin{aligned}
    \det A &=
    \begin{vmatrix}
    1 & 0 & -1 \\
    0 & x & 0 \\
    x & 0 & -1
    \end{vmatrix}
    \\[1em]
    &=
    1
    \begin{vmatrix}
    x & 0 \\
    0 & -1
    \end{vmatrix}
    -
    0
    \begin{vmatrix}
    0 & 0 \\
    x & -1
    \end{vmatrix}
    +
    (-1)
    \begin{vmatrix}
    0 & x \\
    x & 0
    \end{vmatrix}
    \\[1em]
    &=
    1[(x)(-1)-(0)(0)]
    +
    (-1)[(0)(0)-(x)(x)]
    \\[1em]
    &=
    -x + x^2
    \\[1em]
    &=
    x^2-x
    \\[1em]
    &=
    x(x-1)
  \end{aligned}
  \end{equation}

  Para que o determinante seja positivo:
  \begin{equation}
  \begin{aligned}
    x(x-1) &> 0
  \end{aligned}
  \end{equation}

  O produto é positivo quando os fatores têm o mesmo sinal:
  \begin{equation}
  \begin{aligned}
    x<0
    \quad \text{ou} \quad
    x>1
  \end{aligned}
  \end{equation}

  Portanto:
  \begin{equation}
  \begin{aligned}
    x \in (-\infty,0)\cup(1,+\infty)
  \end{aligned}
  \end{equation}
}
